\documentclass[11pt]{article}
\usepackage{manual_docs_style}
\externaldocument{build/run_graph_schema}[run_graph_schema.pdf]
\externaldocument{build/compute_metrics_stage}[compute_metrics_stage.pdf]
\externaldocument{build/generate_question_stage}[generate_question_stage.pdf]

\begin{document}

\docTitle{Dataset Manifest Schema}
\phantomsection\label{docstart:dataset_manifest_schema}

\section*{Overview}

The dataset manifest records which materialized and eligible runs are selected for a concrete benchmark dataset.
It separates the candidate universe from the final train/test samples used to create questions.

The intended flow is:

\begin{DocCode}
resolved run graph
  -> materialized runs
  -> eligibility metrics
  -> selected dataset manifest
  -> exported questions
\end{DocCode}

Candidate generation may intentionally overproduce runs.
The dataset manifest is the frozen selection of examples that a downstream question bundle consumes.

\section*{Location}

Selected dataset artifacts live under:

\begin{DocCode}
models/simulink/<MODEL>/datasets/<dataset_id>/
\end{DocCode}

Recommended files:

\begin{itemize}
\item \code{selected_runs.jsonl}: flat list of selected run IDs and their split/class metadata.
\item \code{sample_manifest.json}: train/test grouping consumed by question generation.
\item \code{selection_report.json}: summary of selection constraints, balance, and rejected candidates.
\end{itemize}

\section*{Inputs}

The selection stage consumes:

\begin{itemize}
\item \code{plans/<policy_id>/run_nodes.jsonl}
\item \code{plans/<policy_id>/run_edges.jsonl}
\item \code{metrics/<policy_id>/eligibility_metrics.jsonl} or \code{runs/eligibility_metrics.json} during migration
\item \code{experiment_config.json}
\item \code{description_levels.json}
\item selection configuration such as \code{all_possible_combinations.csv} or a future \code{selection_policies/<selection_id>.json}
\end{itemize}

\section*{\texorpdfstring{\titlecode{selected_runs.jsonl}}{selected\_runs.jsonl}}

Each line records one selected run.

\begin{DocCode}
{
  "run_id": "run_7fd9daa0055fb3f412ef5e0ab8d4c5ba",
  "family_id": "fam_de4d7c0a5335e9005510cc2442d0971e",
  "split": "test",
  "class_internal": "gravity",
  "class_agent_facing_name": "gravity",
  "policy_id": "benchmark_v1",
  "dataset_id": "paper_main_v1",
  "selection_seed": 123,
  "source": "programmatic",
  "validation_profile": "paper_selected"
}
\end{DocCode}

Allowed \code{split} values are \code{train}, \code{test}, and \code{other}.

\section*{\texorpdfstring{\titlecode{sample_manifest.json}}{sample\_manifest.json}}

\code{sample_manifest.json} groups selected runs into concrete train/test panels for question generation.
It preserves the existing question-generation concept of selecting shared test sets by \code{test_set_slug} and \code{seed}, and train samples by \code{shot_slug}.

Each manifest entry should include:

\begin{itemize}
\item \code{train_samples}
\item \code{train_samples_baselines}
\item \code{test_samples}
\item \code{test_samples_baselines}
\item \code{other_samples}
\item \code{seed}
\item \code{test_set_slug}
\item \code{shot_slug_recipe}
\item \code{train_test_sample_hash}
\end{itemize}

\section*{Selection Rules}

Selection should enforce:

\begin{itemize}
\item only eligible intervention runs may be selected as intervention-class samples;
\item only family-eligible baselines may be selected as no-intervention samples;
\item class counts in each test panel should be as balanced as possible;
\item no test panel should reuse two samples from the same family unless explicitly allowed;
\item train and test samples for the same question must not overlap;
\item selected samples should preserve the shared-test-set pairing required by the experiment plan.
\end{itemize}

\section*{\texorpdfstring{\titlecode{selection_report.json}}{selection\_report.json}}

\code{selection_report.json} records selection diagnostics:

\begin{itemize}
\item timestamp;
\item model;
\item \code{policy_id};
\item \code{dataset_id};
\item number of eligible families;
\item number of selected train, test, and other samples;
\item per-class counts;
\item rejected candidates and rejection reasons;
\item adversarial-baseline checks, when enabled;
\item final pass/fail status for selection acceptance criteria.
\end{itemize}

\section*{Relationship to Question Generation}

Stage: Generate Question consumes the selected dataset manifest and exports the task bundle under:

\begin{DocCode}
tsENV_questions/<MODEL>/
\end{DocCode}

The exported bundle remains the only dependency of agentic rollout.
The dataset manifest is a pre-export selection artifact used to make the candidate-to-question boundary explicit.

\end{document}
